\TNchapter[Agentic systems represent a shift from static automation to adaptive, goal-directed intelligence embedded across the enterprise.]{Quality and Reliability Evaluation }
\begin{TNtwocol}
\section{Response Accuracy and Factual Correctness }
\subsection{Introduction}
Response accuracy forms the foundation of reliable AI agent evaluation. It requires calculating output accuracy, semantic accuracy (meaning preservation), and contextual accuracy across both final responses and reasoning paths.

\subsection{Evaluation Methods}
\begin{enumerate}
    \item \textbf{Exact Match/Structural}: For deterministic tasks or JSON outputs.
    \item \textbf{Semantic Similarity}: Using embeddings to measure distance, though threshold selection is domain-specific.
    \item \textbf{LLM-as-a-Judge}: Using models (e.g., GPT-4o) to provide Likert scores (1-5) and reasoning for Intent Resolution and Completeness.
    \item \textbf{Groundedness}: Ensuring answers are supported by context and declining when information is missing.
\end{enumerate}

\subsection{Best Practices}
Evaluation should be multi-level (intent, tools, context), use configurable thresholds for binary pass/fail, and handle non-determinism via "soft failures" and automatic retries.
\end{TNtwocol}
\begin{TNtwocol}
\section{Hallucination Detection and Measurement}
\subsection{Introduction}
Hallucinations (factual, context, capability, temporal, computational) are critical failure modes where agents generate plausible but incorrect info.

\subsection{Detection Methods}
\begin{itemize}
    \item \textbf{Groundedness Evaluation}: Verifying claims against evidence (e.g., Azure AI Foundry).
    \item \textbf{Multi-Source Verification}: Cross-referencing outputs with trusted sources.
    \item \textbf{Self-Consistency}: Checking consistency across multiple phrasings of the same question.
    \item \textbf{LLM-as-a-Judge}: Using specialized prompts to detect unsupported claims.
\end{itemize}

\subsection{Metrics}
\begin{itemize}
    \item \textbf{Hallucination Rate}: Percentage of responses with hallucinations (Critical < 0.5\%).
    \item \textbf{Groundedness Score}: Average score of how well claims are supported.
\end{itemize}
Mitigation involves anti-hallucination prompting, context expansion, and confidence calibration.
\end{TNtwocol}
\begin{TNtwocol}
\section{Consistency and Reproducibility Testing }
\subsection{Introduction}
Agents introduce non-determinism (model sampling, tool variability). Consistency measures the similarity of outcomes across runs, while Reliability ensures predictable performance.

\subsection{Metrics}
\begin{itemize}
    \item \textbf{Pass\textsuperscript{k} Reliability}: Success rate across $k$ independent runs.
    \item \textbf{Response Variability}: Semantic similarity of outputs across multiple trials.
    \item \textbf{Plan/Tool Consistency}: Stability of reasoning paths and tool selection.
\end{itemize}

\subsection{Best Practices}
Control temperature (lower for consistency), enforce structured output formats (JSON schemas), and use soft failure thresholds to tolerate minor non-determinism.
\end{TNtwocol}
\begin{TNtwocol}
\section{Reliability Metrics}
\subsection{Pass\textsuperscript{k} Metrics}
Pass\textsuperscript{k} evaluates consistency by measuring the percentage of successful completions across $k$ independent trials. It addresses the non-deterministic nature of LLMs where single-shot success is insufficient.

\subsection{Theoretical Foundation}
Based on the $\tau$-bench framework, comparisons against goal states in multi-turn conversations reveal that even advanced models may have low Pass\textsuperscript{8} scores despite decent single-run accuracy.

\subsection{Implementation}
\begin{itemize}
    \item $k=3-5$: Initial development.
    \item $k=8-10$: Pre-production validation.
    \item $k=20+$: High-stakes applications.
\end{TNtwocol}
\begin{TNtwocol}
\end{itemize}
Applications include setting deployment thresholds (e.g., Pass\textsuperscript{8} > 80\%) and detecting model drift.
\section{Edge Case and Boundary Condition Testing }
\subsection{Introduction}
Testing atypical scenarios (Rare Inputs, Extreme Values, Ambiguity) is critical for robustness.

\subsection{Importance}
Edge cases cause trust erosion and cascading failures. Compliance in valid domains (Healthcare, Finance) requires demonstrated robustness.

\subsection{Identification & Testing}
\begin{itemize}
    \item \textbf{Sources}: Production logs, Domain Experts, Red Team exercises.
    \item \textbf{Methodologies}: Golden Prompt Sets (Adversarial, Malformed), Metamorphic Testing (Input transformations), and Fuzzing.
\end{itemize}

\end{TNtwocol}
\begin{TNtwocol}
\subsection{Metrics}
\textbf{Graceful Degradation Score}: Evaluates error clarity and safety preservation when the task cannot be completed.
\section{Robustness Under Varying Inputs }
\subsection{Introduction}
Measures ability to handle input variability (Linguistic, Contextual, Technical) and distributional shifts over time.

\subsection{Evaluation Methodologies}
\begin{itemize}
    \item \textbf{Perturbation-Based Testing}: Character-level (typos), Word-level (synonyms), and Semantic (negation) modifications.
    \item \textbf{Distributional Testing}: Out-of-Distribution (OOD) detection and Cross-Population testing (demographics).
    \item \textbf{Stress Testing}: Volume/Scale stress and Adversarial stress (Prompt Injection).
\end{itemize}
\textbf{Metamorphic Testing} ensures invariance relative to rephrasing or ordering.
\end{TNtwocol}
